\section{Adaptive Mesh Refinement Methodology}
\label{sec:adaptive_refinement}

Adaptive mesh refinement is a well-established strategy for concentrating
computational resolution in selected regions of a numerical domain. Classical
AMR methods provide important historical context for the general concept of
local refinement \cite{berger1984,berger1989}. The cell-selection indicators
developed in the present work are not implementations of those classical AMR
algorithms; they constitute a separate one-shot refinement experiment built
around the OpenFOAM solution and the SOM-derived information.

\subsection{Controlled comparison}

Three cell-selection strategies were evaluated under identical refinement
budgets:

\begin{enumerate}
    \item a conventional velocity-gradient indicator;
    \item direct refinement based on SOM quantization error (SOM-QE);
    \item a hybrid indicator combining the velocity gradient with
    SOM-derived regime-transition information.
\end{enumerate}

For each method, refinement budgets of 5\%, 10\%, 15\%, and 20\% of the
1600-cell M40 baseline mesh were considered.

The corresponding numbers of initially selected cells were therefore 80,
160, 240, and 320 cells, respectively.

The same local refinement operation was applied after cell selection.
Consequently, the final adaptive meshes contained 1840, 2080, 2320, and
2560 cells for the four refinement budgets.

This design isolates the effect of the cell-selection criterion: at a given
budget, each method starts from the same M40 solution, selects the same
number of cells, and uses the same refinement procedure.

\subsection{Velocity-gradient indicator}

The conventional baseline ranks cells according to the local velocity-gradient
magnitude,

\begin{equation}
G_i = |\nabla\mathbf{U}_i|.
\end{equation}

Cells with the largest values of \(G_i\) are selected until the prescribed
refinement budget is reached.

This method provides a physically interpretable conventional baseline:
regions in which the velocity field varies rapidly receive higher refinement
priority.

\subsection{Direct SOM quantization-error indicator}

The first SOM-based hypothesis uses the cell-level quantization error directly
as a refinement indicator,

\begin{equation}
I_{\mathrm{QE},i}=QE_i.
\end{equation}

Cells with the largest quantization errors are selected first.

This criterion tests whether CFD states that are poorly represented by the
trained SOM also correspond to locations where additional spatial resolution
is beneficial.

No fine-grid reference information is used during this selection.

\subsection{SOM regime-transition indicator}

The BMU classification provides a discrete label \(b_i\) for every M40 cell.
To identify interfaces between different SOM-classified flow states, a local
transition indicator is defined using the direct structured-grid neighbors of
cell \(i\):

\begin{equation}
T_i =
\frac{
\displaystyle
\sum_{j\in\mathcal{N}_i}
\mathbb{I}(b_j\neq b_i)
}{
|\mathcal{N}_i|
},
\end{equation}

where \(\mathcal{N}_i\) is the set of existing direct neighbors and
\(\mathbb{I}\) is the indicator function.

Thus,

\begin{equation}
0\leq T_i\leq1.
\end{equation}

A value of zero indicates that all direct neighbors belong to the same SOM
state, while larger values indicate increasing local interaction between
different SOM-classified states.

The transition indicator is constructed exclusively from the M40 SOM
classification. It does not use the M160 reference solution.

The cell-level correlation between the transition indicator and the
velocity-gradient magnitude was approximately

\begin{equation}
r(G,T)=0.261,
\end{equation}

substantially smaller than the correlation between gradient magnitude and
quantization error,

\begin{equation}
r(G,QE)=0.879.
\end{equation}

The transition indicator therefore contains information that is less
redundant with the conventional gradient criterion than the direct
quantization error.

\subsection{Hybrid-v1 indicator}

The hybrid indicator was designed to preserve the conventional gradient as
the primary refinement signal while allowing SOM regime-transition
information to modify the ranking of cells.

The normalized gradient is defined as

\begin{equation}
G_i^{*} =
\frac{G_i}{\max_j G_j}.
\end{equation}

The Hybrid-v1 indicator is then

\begin{equation}
I_{\mathrm{hyb},i}
=
G_i^{*}
\left(1+\alpha T_i\right),
\end{equation}

with

\begin{equation}
\alpha=0.5.
\end{equation}

The coefficient was fixed for the Hybrid-v1 experiment and was not optimized
against the M160 reference.

This distinction is important because tuning the coefficient to minimize
M160 discrepancies and subsequently evaluating performance against the same
M160 solution would contaminate the validation comparison.

Hybrid-v1 should therefore be interpreted as a predefined exploratory
indicator rather than an optimized model.

\subsection{Selection overlap}

Because the gradient remains the dominant component of Hybrid-v1, the hybrid
criterion modifies only a relatively small fraction of the conventional
gradient selections.

The overlap between Gradient and Hybrid-v1 selections was

\begin{table}[htbp]
\centering
\caption{Overlap between Gradient and Hybrid-v1 cell selections.}
\label{tab:gradient_hybrid_overlap}
\begin{tabular}{rrrrr}
\toprule
Budget & Selected & Common & Exclusive per method & Overlap \\
\midrule
5\%  & 80  & 73  & 7  & 91.25\% \\
10\% & 160 & 149 & 11 & 93.12\% \\
15\% & 240 & 227 & 13 & 94.58\% \\
20\% & 320 & 306 & 14 & 95.62\% \\
\bottomrule
\end{tabular}
\end{table}

Consequently, Hybrid-v1 is not a replacement of the gradient criterion by a
machine-learning criterion. It is a controlled perturbation of the
gradient-based ranking using information extracted from SOM flow-state
boundaries.


\subsection{Second-stage SOM-interface refinement criterion}

The results obtained with direct SOM quantization error motivated a change in
the evaluation question. Rather than treating the SOM as a direct estimator
of discretization error, the subsequent analysis asked whether the SOM
classification contains information about \emph{where refinement is likely to
be useful}.

The BMU field provides a spatial classification of local flow states. Cells
located at interfaces between different BMU states were therefore considered
potentially informative because they represent locations where the local flow
regime changes in the SOM feature space.

For an interface cell, the local transition fraction \(T_i\) was combined
with two quantities already available from the M40 pilot solution. The first
was the normalized velocity-gradient magnitude,

\begin{equation}
G_i^* =
\frac{G_i}{\max_j G_j},
\end{equation}

and the second was the normalized maximum velocity-vector jump across a
direct SOM interface,

\begin{equation}
\Delta U_i^* =
\frac{
\max_{j\in\mathcal{N}_i,\;b_j\neq b_i}
\|\mathbf{U}_i-\mathbf{U}_j\|_2
}{
\max_k \Delta U_k
}.
\end{equation}

The resulting SOM-interface ranking was defined as

\begin{equation}
I_{\mathrm{SI},i}
=
T_i
\sqrt{G_i^*\,\Delta U_i^*}.
\end{equation}

Only cells for which \(T_i>0\) were eligible for interface selection. The
highest-ranked 40 interface cells were then frozen as common seeds. Starting
from exactly those same seeds, the selected refinement region was expanded by
a fixed graph distance: one face-sharing layer for the \(1N\) case and two
layers for the \(2N\) case.

This construction deliberately separates two questions. The SOM determines
\emph{where flow-state transitions occur}, while the gradient and
velocity-jump terms determine how strongly those transitions are ranked.
Neither the M160 solution nor any reference-derived error field is used in
this construction.

The resulting experiments selected 78 cells for \(1N\) and 114 cells for
\(2N\). These counts were not forced to equal a nominal 10\% budget.
Instead, they were subsequently reproduced exactly with conventional
Gradient selection, creating cell-matched controls with identical final
cell counts. This comparison isolates the spatial-selection mechanism while
avoiding a confounding difference in total mesh size.

\subsection{Evaluation against the fine-grid reference}

After refinement, each adaptive case is solved independently with OpenFOAM.
The resulting velocity field is compared with the M160 fine-grid numerical
reference.

Because the adaptive and reference meshes contain different cell-center
locations, the M160 solution is interpolated to the adaptive-mesh cell
centers before the velocity discrepancy is evaluated.

For cell \(i\), the nondimensional velocity-vector discrepancy is

\begin{equation}
e_i =
\frac{
\left\|
\mathbf{U}_{i,\mathrm{adaptive}}
-
\mathbf{U}_{i,\mathrm{M160}}
\right\|_2
}{
U_{\mathrm{lid}}
}.
\end{equation}

Since adaptive cells have different volumes, global mean and RMS measures are
volume weighted.

The M160 solution is used only at this evaluation stage and not during the
selection of cells for Gradient, SOM-QE, or Hybrid-v1.
